%%%%%%%%%%%%%%%%%%%%%%%
%%% DOCUMENT SETTINGS
%%%%%%%%%%%%%%%%%%%%%%%
\documentclass[11pt]{exam}
\printanswers % Uncomment to show answers if desired
%%%%%%%%%%%%%%
%%% PACKAGES
%%%%%%%%%%%%%%
\usepackage{amsmath,amsfonts,amssymb,amsthm}  % math fonts and symbols
\usepackage{bbm} % Math blackboard font
\usepackage{enumitem} % Custom list environments
\usepackage{textcomp} % Misc symbols
\usepackage[top=2cm, bottom=2cm, right=1.2cm, left=1.2cm, paper=a4paper, headheight=6cm]{geometry} % Page layout
\usepackage{xcolor}  % Colors
\usepackage{graphicx} % Graphics
%----------------------------------- TikZ
\usepackage{tikz}
\usepackage{pgfplots}
\pgfplotsset{compat=1.17}
\usepgfplotslibrary{fillbetween}
\usetikzlibrary{shapes,arrows,matrix,positioning,calc}
\newcommand{\tikzmark}[1]{\tikz[baseline,remember picture] \coordinate (#1) {};}
%%%%%%%%%%%%%%%%%%%%%%%%%%
%%% COLORS
%%%%%%%%%%%%%%%%%%%%%%%%%%
\definecolor{brickred}{RGB}{200, 50, 50}
\definecolor{skyblue}{RGB}{0, 150, 220}
\definecolor{olivegreen}{RGB}{120, 160, 60}
%%%%%%%%%%%%%%%%%%%%%%%%%%
%%% MACROS
%%%%%%%%%%%%%%%%%%%%%%%%%%
\newcommand{\N}{\mathbb{N}} % natural numbers
\newcommand{\R}{\mathbb{R}} % real numbers
\makeatletter
\renewcommand{\@makefnmark}{\textsuperscript{\arabic{footnote}}}
\renewcommand{\thefootnote}{\arabic{footnote}}
\makeatother
\setcounter{footnote}{0}
\newcommand{\BAR}{%
  \hspace{-\arraycolsep}\strut\vrule\hspace{-\arraycolsep}}
%%%%%%%%%%%%%%%%%%%%%%%%%%%%%%%%%
%%% HEADER AND FOOTER
%%%%%%%%%%%%%%%%%%%%%%%%%%%%%%%%%
% Document information
\newcommand{\degree}{Bachelor in Philosophy, Politics and Economics}
\newcommand{\shortdeg}{PPE}
\newcommand{\exam}{Mock Exam 1}
\newcommand{\subject}{Mathematics}
\newcommand{\theyear}{2024/25}
\newcommand{\ccauthor}{A. G. Azze}
\newcommand{\thedate}{May 4, 2025}
\newcommand{\university}{CUNEF Universidad}
% Set header and footer
\pagestyle{headandfoot}
\header{\degree \\ \subject\ - \theyear}{}{\thedate \\ \ccauthor}
\footer{\university}{}{\thepage}
%%%%%%%%%%%%%%%%%%%%%%%%%%
%%% DOCUMENT CONTENT
%%%%%%%%%%%%%%%%%%%%%%%%%%
\begin{document}
%--- Instructions ---%
\begin{center}
    {\Huge \textbf{\exam{}}}
    \begin{flushright}
        \vspace{-0.3cm}
        \textbf{Time Limit:} 120 minutes
    \end{flushright}
    \vspace{-0.1cm}
    \fbox{\parbox{\textwidth}{\centering 
        \begin{enumerate}[leftmargin=0.5cm, label=$\bullet$, rightmargin=0.2cm]
            \item Every non-trivial calculation and claim must be justified.
            % \item Calculators are not necessary and will not be allowed.
            \item Digital devices and internet access are prohibited.
        \end{enumerate}
    }}
\end{center}

\begin{questions}

%%%%%%%%%%%%%%%%%%%%%%%%%%%%%%%%%%%%%%%%%%%%%%%%%%%%%%%%%%%%%%%%%%%%%
% Question 1: Lorenz Curves & Gini
%%%%%%%%%%%%%%%%%%%%%%%%%%%%%%%%%%%%%%%%%%%%%%%%%%%%%%%%%%%%%%%%%%%%%
\question[3]
In the city of Greenvale, community gardens share irrigation hours among volunteer groups.  The Lorenz curve for weekly irrigation-hours is
\[
L(x)=\frac{3}{5}x^p + a\,x,\quad x\in[0,1],
\]
where \(L(x)\) is the fraction of total hours used by the least‑served \(100x\)% of groups.
\begin{enumerate}[leftmargin=0.5cm,label=(\alph*),rightmargin=0.2cm]
  \item Show that \(L(1)=1\) implies \(a=\tfrac{2}{5}\).
  \item Compute the area under the Lorenz curve, \(B=\int_0^1L(x)\,dx\), in terms of \(p > 0\).
  \item Given that the Gini index satisfies \(G=1-2B\) and is measured to be \(0.30\), find \(p\).
\end{enumerate}

\begin{solution}
\begin{enumerate}[label=(\alph*)]
  \item From \(L(1)=\tfrac35 + a=1\), so \(a=\tfrac25\).
  \item 
  \[ 
    B=\int_0^1\Bigl(\tfrac35x^p + \tfrac25x\Bigr)dx
    =\frac{3}{5(p+1)} + \frac{1}{5}.
  \]
  \item Since \(G=1-2B=0.30\),
  \[
    1-2\Bigl(\frac{3}{5(p+1)} + \frac{1}{5}\Bigr)=0.30
    \;\Longrightarrow\;
    p=3.
  \]
\end{enumerate}
\end{solution}

%%%%%%%%%%%%%%%%%%%%%%%%%%%%%%%%%%%%%%%%%%%%%%%%%%%%%%%%%%%%%%%%%%%%%
% Question 2: Input–Output
%%%%%%%%%%%%%%%%%%%%%%%%%%%%%%%%%%%%%%%%%%%%%%%%%%%%%%%%%%%%%%%%%%%%%
\question[3]
A region has two green industries: Waste Management (W) and Recycling (R).  To produce one tonne of output:
\[
\text{W needs }0.3\text{ t of W and }0.2\text{ t of R},\quad
\text{R needs }0.4\text{ t of W and }0.3\text{ t of R}.
\]
Final demand is \(d_W=50\) t and \(d_R=30\) t.  Let outputs be \(x_W,x_R\).
\begin{enumerate}[leftmargin=0.5cm,label=(\alph*),rightmargin=0.2cm]
  \item Write the equilibrium equations “output = intermediate + final.”
  \item Cast them in matrix form \(A\,x=d\).
  \item Compute the determinant of \(A\) and conclude whether the system has a unique solution.
  \item Solve for \((x_W,x_R)\).
\end{enumerate}

\begin{solution}
\begin{enumerate}[label=(\alph*)]
  \item Using \(a_{ij}\) = tonnes of input \(i\) per tonne of output \(j\),
  \[
    x_W = 0.3x_W + 0.4x_R + 50,\quad
    x_R = 0.2x_W + 0.3x_R + 30.
  \]
  \item Rearranged:
  \[
    \begin{pmatrix}
      0.7 & -0.4\\
     -0.2 & 0.7
    \end{pmatrix}
    \begin{pmatrix}x_W\\x_R\end{pmatrix}
    =
    \begin{pmatrix}50\\30\end{pmatrix}.
  \]
  \item \(\displaystyle \det A = 0.7\cdot0.7 - (-0.4)(-0.2) = 0.49-0.08 = 0.41\neq 0\), so a unique solution exists.
  \item For instance, multiply the equations by 10:
  \[
    7x_W - 4x_R = 500,\qquad -2x_W + 7x_R = 300.
  \]
  Solving this system gives
  \[
    x_R = \frac{3100}{41} \approx 75.61,\qquad
    x_W = \frac{4700}{41} \approx 114.63.
  \]
\end{enumerate}
\end{solution}

%%%%%%%%%%%%%%%%%%%%%%%%%%%%%%%%%%%%%%%%%%%%%%%%%%%%%%%%%%%%%%%%%%%%%
% Question 3: Limits & Continuity
%%%%%%%%%%%%%%%%%%%%%%%%%%%%%%%%%%%%%%%%%%%%%%%%%%%%%%%%%%%%%%%%%%%%%
\question[2]
In a wildlife reserve, the density of a rare flower at distance \(x\) (m) from water is
\[
f(x)=\frac{\sqrt{3x+1}-2}{x-1},\quad x\neq1.
\]
\begin{enumerate}[leftmargin=0.5cm,label=(\alph*),rightmargin=0.2cm]
  \item Determine the domain and where \(f\) is continuous.
  \item Compute \(\displaystyle\lim_{x\to1}f(x)\).
  \item Find a value for \(f(1)\) that makes \(f\) continuous at \(x=1\).
\end{enumerate}

\begin{solution}
\begin{enumerate}[label=(\alph*)]
  \item Domain: \(x\ge-1/3\), \(x\neq1\); continuous on its domain.
  \item By L’Hôpital, limit = \(3/4\).
  \item Set \(f(1)=3/4\) to restore continuity.
\end{enumerate}
\end{solution}

%%%%%%%%%%%%%%%%%%%%%%%%%%%%%%%%%%%%%%%%%%%%%%%%%%%%%%%%%%%%%%%%%%%%%
% Question 4: Improper Integrals
%%%%%%%%%%%%%%%%%%%%%%%%%%%%%%%%%%%%%%%%%%%%%%%%%%%%%%%%%%%%%%%%%%%%%
\question[2]
Determine convergence and compute where possible:
\begin{enumerate}[label=(\alph*)]
\item \(\displaystyle\int_{0}^{\infty}e^{-2x}\,dx\).
\item For which \(p>0\) does \(\int_{1}^{\infty}\frac{1}{x^p}\,dx\) converge?
\item Using comparison of the functions $1/\sqrt{x}$ and $1/x$, decide whether \(\int_{1}^{\infty}\frac{1}{\sqrt{x}}\,dx\) converges or diverges.
\item \(\displaystyle\int_{0}^{\infty}x e^{-x}\,dx\). Use the fact that, for $F(x) = - (x+1)e^{-x}$, $F'(x) = x e^{-x}$.
\end{enumerate}

\begin{solution}
\begin{enumerate}[label=(\alph*)]
\item \(\int_0^\infty e^{-2x}dx = -\tfrac12 (e^{-2x})_0^\infty = 1/2\).
\item Converges iff \(p>1\).
\item Since \(1/\sqrt{x}=x^{-1/2}\) and \(-1/2\le-1\)? Actually compare to \(1/x\): for large x, \(1/\sqrt{x}>1/x\), and \(\int1/x\) diverges, so this diverges.
\item Use \(\int x e^{-x}dx = - (x+1)e^{-x}\).  At \(0\) gives \(-(0+1)\cdot1=-1\); at \(\infty\to0\).  Difference = \(1\).
\end{enumerate}
\end{solution}

\end{questions}
\end{document}
